\documentclass[11pt,a4paper]{article}
% Compile with XeLaTeX (fontspec and xeCJK).
\usepackage{fontspec}
\setmainfont{LibertinusSerif}[Extension=.otf,UprightFont=*-Regular,BoldFont=*-Bold,ItalicFont=*-Italic,BoldItalicFont=*-BoldItalic]
\setsansfont{LibertinusSans}[Extension=.otf,UprightFont=*-Regular,BoldFont=*-Bold,ItalicFont=*-Italic]
\usepackage{xeCJK}
\setCJKmainfont{Hiragino Mincho ProN}
\xeCJKDeclareCharClass{Default}{"00B0, "00B7, "2013, "2014, "2018 -> "201F, "2026}
\usepackage[english]{babel}
\usepackage[a4paper,margin=2.5cm,top=2.3cm]{geometry}
\usepackage{graphicx}
\usepackage{xcolor}
\usepackage{titlesec}
\usepackage{enumitem}
\usepackage{fancyhdr}
\usepackage{enumitem}
\usepackage{xurl}
\usepackage[hidelinks]{hyperref}
\definecolor{accent}{RGB}{31,73,125}
\definecolor{muted}{RGB}{93,104,119}
\hypersetup{colorlinks=true,urlcolor=accent,linkcolor=accent,pdfauthor={Romain Fernex},pdftitle={Report on the EU–South Korea Free Trade Agreement: a discussion on the history and future of European trade policy in Asia}}
\urlstyle{same}
\setlist[itemize]{topsep=.25em,itemsep=.15em,parsep=0pt,leftmargin=1.4em}
\setcounter{secnumdepth}{0}
\titleformat{\section}{\large\bfseries\color{accent}}{}{0em}{}
\titlespacing*{\section}{0pt}{1.6em}{.6em}
\setlength{\parindent}{1.2em}
\setlength{\parskip}{0.35em}
\pagestyle{fancy}
\fancyhf{}
\fancyhead[L]{\small\color{muted}Romain Fernex}
\fancyhead[R]{\small\color{muted}\itshape EU–South Korea Free Trade Agreement}
\fancyfoot[C]{\small\color{muted}\thepage}
\renewcommand{\headrulewidth}{0pt}
\newcommand{\partbreak}{\par\begin{center}\color{muted}$\ast$\quad$\ast$\quad$\ast$\end{center}\par}


\begin{document}
\thispagestyle{empty}
\noindent{\small\sffamily\color{accent}\MakeUppercase{Report · European trade policy}}\par\medskip
\noindent{\LARGE\bfseries Report on the EU–South Korea Free Trade Agreement\par}
\smallskip\noindent{\large\itshape A discussion on the history and future of European trade policy in Asia\par}
\medskip
\noindent{\color{accent}\rule{\linewidth}{0.8pt}}\par\smallskip
\noindent{\small\color{muted}Romain Fernex · 2023 · Report in English}\par
\smallskip\noindent{\small\color{muted}\emph{Article also published in French on:} \url{https://www.callonu-escpeurope.com/articles/l-accord-de-libre-échange-ue-corée-du-sud-discussion-sur-l-histoire-et-l-avenir-de-la-politique-commerciale-européenne-en-asie}}\par\bigskip

On 2 November 2022, William Yuen Yee, research assistant for the Columbia-Harvard China and the World program, published an article for The Diplomat on the great difficulties encountered by the EU in signing a new free trade agreement (FTA) with Indonesia. It points to the setbacks of an EU that is quick to assert its desire to expand its economic relations with Asia{[}1{]} in an increasingly tense international context, particularly with the growing number of supply chain disruptions, but which is often held back by its will to impose environmental standards on its partners. Standards which are seen as insurmountable constraints by the vast majority of states in the Indo-Pacific region. Generally speaking, Europe’s economic presence in Asia was relatively late and marked from the outset by a notable ambition to defend strict social and environmental standards. In order to better understand this specificity and the debates that ensue, this article offers to look back at the EU’s very first free trade agreement in Asia, the EU-Korea FTA. A deal that was a precursor to Europe’s current trade policy in Asia and more generally to its conception of international economic relations at the bilateral level. While it is impossible to summarize a decade of EU-Asia trade in one article, we will try to see how the FTA signed in 2007 between the European Union and South Korea reflects the particularity and ambiguity of the EU’s economic relations with Asia.

To this end, we will first look at the composition of the agreement itself, focusing on those chapters that are most innovative. Then we shall discuss the economic impact of the agreement on relations between South Korea and Europe. Finally, this article will broach the various controversies surrounding the implementation of the agreement and extend our reflection to some of the FTAs signed by the EU in Asia that followed the KOREU FTA.

\partbreak

As in all FTAs, several clauses in Chapter 2 of the agreement aim at the almost total elimination of existing tariff barriers – in the order of 99\% for industrial products – with the exception of certain specific sectors such as fisheries or agriculture. More interestingly, this chapter also contains provisions concerning non-tariff barriers, i.e. obstacles to trade other than customs duties {[}2{]}, such as the existence of certain country-specific standards. For example, in the automotive sector, the agreement provides for the adoption of international safety standards by South Korea in order to harmonize standards between the two countries. A “fast-track dispute settlement” system is also put in place to allow for quicker arbitration in case of a dispute. This attention to non-tariff barriers is reflected in Chapter 4 on technical barriers to trade and Chapter 6 on customs and trade facilitation. Finally, Chapter 10 on intellectual property is of real interest to the European Union, much more so than to Korea, and aims to guarantee, among other things, respect for indications linked to the geographical origin of certain products, particularly in the food industry. Thus PDO-style designations etc. benefit from a higher level of protection. In the same spirit, the agreement includes a “rule of origin” principle which imposed a minimum threshold of 45\% of components manufactured in the country of origin for the final exported product to be considered as originating there.

Beyond these mainly commercial considerations, one chapter is particularly innovative: Chapter 13 entitled Trade and Sustainable Development. Like the agreements put in place by the EU for the integration of Eastern European countries, the signing of the agreement is conditional on the acceptance of a significant number of environmental and social standards. Institutional structures are set up by each party in the FTA, with the presence of civil society actors, to ensure that these commitments are respected. Thus, as regards labor law, the two parties undertake to scrupulously follow the rules laid down by the International Labor Organization (ILO), even though South Korea had only ratified 4 of the 8 major ILO conventions {[}3{]}. On the environmental front, the two parties promise, among other things, to put in place concrete means to meet their commitments to the various environmental agreements they have signed – within the framework of the COPs for example.

This set of measures, although still limited, was a first for the time and the free trade agreement negotiated by the US with South Korea at the same time did not include any such provision.

\partbreak

But what about the economic and social impact of this agreement for both parties? For the European Union, whose total export volume to Korea increased by 70\% from 2010 to 2020, the deal is a success. An increase that does not only concern manufactured goods, but also service exports are these also rose by 50\% over the same period. This increase benefits several key EU industries such as automotive and electronics. Moreover, several surveys carried out by EEAS have highlighted the positive reception of this agreement by European companies present in Korea. Three quarters of them mentioned that they have expanded their activity in the ROC since the agreement and show a satisfaction rate close to 90\%.

However, on the Korean side, the enthusiasm preceding the agreement quickly gave way to a certain disenchantment in the face of disappointing performances for Korean exports. In fact, over the same period, Korean exports to the EU increased by only 32\%{[}4{]}, while imports from Europe rose much more rapidly, particularly in the five years following the signing of the agreement. This new trend even leads to a reversal of the trade balance between the two sides in favor of the EU from 2011 to 2017, which surprises the Koreans who were previously in net excess. The reasons put forward for this difference in performance are numerous: From a cyclical point of view, the EU, even more so than Korea, suffered from the impact of the 2008 financial crisis and the European domestic market proved to be much less dynamic than expected. Besides, Korean manufactured goods exports to the EU also declined. Indeed, Korean industries exporting to Europe are reorganizing to relocate production to Europe where previously they were limited to distribution. Thus, Korean FDIs{[}5{]} towards Eastern European countries such as the Czech Republic increased and large Korean companies such as Kia Motors set up their first factories in Europe to offshore their production for the European domestic market. Finally, in addition to this relocation movement in Eastern Europe, Korean companies started relocating more and more to South-East Asia, particularly to Vietnam. In fact, South Korea is the leading country for FDI to Vietnam, with more than 19 billion dollars invested since 2008. This trend is largely driven by the Korean industrial giants, the chaebolgs, of which Samsung is one of the largest. Therefore, a significant part of Korean products are exported to Vietnam for assembly before being re-exported to Europe, which means that these exports are not included in the EU-Korea balance.

Therefore, from a trade perspective this agreement remains positive overall and performs respectably, although disappointingly, compared to Korean expectations. Nevertheless, the controversy surrounding this agreement is not so much focused on this commercial aspect as on other aspects that are more unique to European trade policy.

\partbreak

Firstly, the EU’s insistence on imposing strong social and environmental standards on its partners has been poorly received by Koreans. Several of the strong principles of the ILO Charter that had not been ratified before the agreement are still not ratified today. For instance, the article on freedom of association, which aims, among other things, to facilitate the establishment of trade unions, is still widely unpopular among Korean Businesses. On 5 July 2019 The EU therefore called for a panel of experts to be set up to investigate possible violations of chapter 13 of the agreement by South Korea. The aforementioned pundits then ruled in favor of the EU’s accusation, which then threatened South Korea with possible trade sanctions. Still in the field of human rights, the legislation around the death penalty is also an important topic of debate. Although the agreement did not explicitly call for its abolition, this major cultural and judicial difference remains important for understanding some of the tensions that remain between the EU and South Korea and which indirectly impact their trade relations. Europe has in fact called on Korea several times to abolish the death penalty, an interference that is sometimes badly received by part of Korean public opinion. A similar tension can be found in the context of negotiations between the EU and Indonesia, with a more environmental than social perspective. Indeed, the European ban on palm oil makes Indonesia fear an asymmetrical relationship at the end of any free-trade agreement with EU and is therefore more reluctant to commit to such a deal compared to other less environmentally demanding agreements such as the one it has signed with its ASEAN partners.

Moreover, technical and economic barriers still exist between the two countries, particularly around the “Rule of Origin”. Indeed, while the EU produces the majority of its export goods within the limits of the European domestic market (notably by relocating to Eastern Europe for the richest countries), Korea often resorts to relocation in foreign countries. Moreover, certain phytosanitary standards are sometimes debated, such as the one concerning the pasteurisation of cheese, a process that is compulsory in Korea but not always used in Europe to the benefit of other production techniques, such as the one used for Parmesan cheese, which was only authorised for export well after the treaty. Finally, the trade deficit encountered by Korea leaves a very mixed memory of the agreement in Korean public opinion. For example, in 2016, at a press conference organised by the EU to celebrate the 5th anniversary of the EU-Korea FTA, no Korean government representatives were present.

In conclusion, as evidenced by the EU’s FTA with Korea and its recent failures with Indonesia, the EU’s desire to establish itself as a “normative power”{[}7{]} and a commercial power sometimes comes up against cultural and economic realities that are antagonistic to the achievement of its objectives. However, this is also what makes the European Union special: its aspiration to spread, including through trade, a mode of development that it considers more sustainable. Although it is still difficult to estimate precisely to what extent this approach hinders the defence of its interests in Asia, it is certain that this approach remains one of the keys of the EU’s differentiation strategy in an area that is economically and ideologically dominated by the United States and China.

\textbf{Notes:}

\textbf{{[}1{]}} Indeed, in an EEAS report of 21 February 2022, the EU presented its plan to “accelerate its strategic engagement in the critically important Indo-Pacific region”. This engagement is part of a dynamic of diversification of the EU’s trading partners. Indeed, the EU announced in early June 2022 that it would step up its efforts to sign new FTAs with other countries in order to reduce the dependency problem it faces in the context of the war between Ukraine and Russia.
\begin{itemize}
  \item Sources:
  \begin{itemize}
    \item EU fact sheet on the EU strategy for cooperation the Indo-pacific, EEAS, 21 February 2022
    \item Philip Blenkinsop, EU seeks to speed up trade deals as Ukraine war forces rethink, 3 June 2022, Reuters
  \end{itemize}
\end{itemize}

\textbf{{[}2{]}} Non-tariff barriers include, for example, embargoes, licensing and quota policies.

\textbf{{[}3{]}} The 4 ILO axes that Korea has not yet ratified are the following:
\begin{itemize}
  \item Freedom of association
  \item Right to organise and collective bargaining
  \item Prohibition and abolition of forced labour
\end{itemize}

\textbf{{[}4{]}} This figure corresponds to the evolution of Korea’s exports to the EU from 2010 to 2020. In reality, the year after the agreement, exports to the EU first decreased and then increased again, especially in 2017. A decline partly offset by the anticipation effect of the agreement which caused a significant increase in exports before its entry into force

\textbf{{[}5{]}} FDI or direct investment abroad corresponds to “investments that a resident institutional unit of an economy makes with the aim of acquiring a lasting interest in a resident institutional unit of another economy and exerting a significant influence on its management within the framework of a long-term relationship”, definition given by INSEE, on 05/11/2019

\textbf{{[}6{]}} According to data published by Statista in April 2022, Indonesia is the world’s largest producer of palm oil and accounts for up to 60\% of global production.
\begin{itemize}
  \item Link to Statista report: \url{https://fr.statista.com/infographie/27346/evolution-de-la-production-mondiale-huile-de-palme-et-principaux-pays-producteurs-en-2022/}
\end{itemize}

\textbf{{[}7{]}} I am using the term used by Zaki Laïdi, a French political scientist and director of research at the CEVIPOF, in his book EU foreign policy in a Globalized world, which contrasts it with more traditional conceptions of power such as American ‘hard power’.

\section{Sources}
\small

\begin{itemize}[leftmargin=1.3em,itemsep=.18em,topsep=.25em]
  \item William Yuen Yee, Can the EU and Indonesia Sign their elusive free trade agreement?, The Diplomat, 2 November 2022.
  \item Alexandra Hennessy and Poppy S. Winanti, Part IV, Chapter 14, EU-Indonesia Trade relations, A geo-economic turn in Trade policy? EU Trade Agreements in the Asia- Pacific Region, 2022, Palgrave Macmillan
  \item EU Strategy for Cooperation in the Indo-Pacific, European Union External Action, 21 February 2022
  \item Evaluation of the results of the EU-Korea Free Trade Agreement, European Union External Action, 2016
  \item Csilla Lakatos and Lars Nilsson, The EU-Korea Free Trade Agreement: anticipation, trade policy uncertainty and impact, note by the EU Chief Economist , June 2015.
  \item Dong-Hee Joe, A review of the First Decade of the Korea-EU FTA, World Economy Brief, vol. 12 No. 27, KIEP, 4 July 2022.
  \item Sunghoon Park, The new politics of trade negotiations: the case of the EU- Korea FTA, Journal of European Integration, 29 November 2017.
  \item Yoo-Duk Kang, Korea-Eu FTA: breaking new ground, Australia, the European Union and the New Trade agenda, 2017, ANU Press.
  \item Samuel Wejchert, EU-South Korea FTA: A ten-year perspective, European Institute for Asian studies, policy brief of 8 September 2021.
  \item European Commission website, section on EU trade relations with South Korea: facts, figures and latest developments.
  \item European Commission fact sheet on the EU-South Korea Free Trade Agreement: a quick reading guide
  \item European Commission fact sheet on the EU-South Korea Free Trade Agreement: 10 key benefits for the European Union
  \item Yoo-Duk Kang and Jun Yeup Kim, Trade and investment between Korea and EU after the Korea-EU FTA and its prospect, KIEP World economy update, vol.3 no.44, 27 September 2013.
  \item Hanns Günther Hilpert and Sunghoon Park, EU and South Korea: a model example for foreign trade policy relations, Part IV, Chapter 10, A geo-economic turn in Trade policy? The EU’s trade agreements in the Asia-Pacific region, 2022, Palgrave Macmillan.
  \item Zaki Laidi, EU foreign policy in a globalised world, normative power and social preferences, 2008.
  \item Trường Lăng, The Journey of Samsung in Vietnam, Viettonkin consulting, 26 July 2022.
  \item Jeon Jong-Hwi, South Korea to become a global test case in its trade dispute with the EU, Hankyoreh, 5 July 2019.
\end{itemize}

\end{document}
